% =============================================================================
% P3A -- Lições aprendidas (PT-BR)
% Formatação: ABNT (classe abnTeX2, modo article), A4, Latin Modern, 12 pt,
% margens 3/2/3/2 cm, espaçamento 1,5, recuo de parágrafo 1,25 cm.
% Figuras esperadas em out/ (compiladas a partir de figuras/):
%   Fig1..Fig5_PT.pdf
% ARQUIVO GERADO por ferramentas/gerar_tex.py a partir de ../MD. Para mudar o
% conteúdo, edite o Markdown e gere de novo, ou edite este .tex diretamente.
% =============================================================================
\documentclass[12pt,a4paper,article,oneside,english,brazil]{abntex2}

\usepackage[T1]{fontenc}
\usepackage[utf8]{inputenc}
\usepackage{lmodern}           % Latin Modern (fonte padrão do LaTeX, com T1)
\usepackage{amsmath}
\usepackage{graphicx}
\usepackage{array}
\usepackage{tabularx}
\usepackage{longtable}
\usepackage[expansion=false]{microtype}
\usepackage{url}
\Urlmuskip=0mu plus 1mu\relax

% Margens ABNT: superior 3 cm, inferior 2 cm, esquerda 3 cm, direita 2 cm
\setlrmarginsandblock{3cm}{2cm}{*}
\setulmarginsandblock{3cm}{2cm}{*}
\checkandfixthelayout

\setlength{\parindent}{1.25cm}
\setlength{\parskip}{0pt}
\setlength{\emergencystretch}{3em}
\setcounter{secnumdepth}{3}
\setcounter{tocdepth}{3}
\renewcommand{\contentsname}{SUMÁRIO}
\renewcommand{\listfigurename}{LISTA DE ILUSTRAÇÕES}
\renewcommand{\listtablename}{LISTA DE TABELAS}

% Tabelas: sem filetes verticais (padrão IBGE adotado pela ABNT)
\renewcommand{\arraystretch}{1.2}

% Fonte abaixo de tabelas e figuras (ABNT NBR 14724)
\newcommand{\fontetab}[1]{\par\vspace{3pt}{\small\noindent Fonte: #1\par}}

% Citações numéricas: o número de cada referência é definido uma vez (ordem
% alfabética da lista) e usado por \refnum{chave}, \cit{chave} e \refitem.
\newcommand{\defref}[2]{\expandafter\gdef\csname refnum@#1\endcsname{#2}}
\newcommand{\refnum}[1]{\csname refnum@#1\endcsname}
\newcommand{\cit}[1]{[\refnum{#1}]}
\newcommand{\refitem}[2]{\par\noindent\hangindent=1.2cm\hangafter=1\makebox[1.2cm][l]{[\refnum{#1}]}#2\par}


% Integrantes do grupo (preencher; a capa usa esta macro)
\newcommand{\integrantes}{[nomes dos integrantes do grupo]}

\begin{document}
\SingleSpacing
\thispagestyle{plain}

\begin{center}\textbf{RELATO DE LIÇÕES APRENDIDAS}\end{center}

\textbf{Grupo SIPA Privacidade 2}

\subsection*{O que aprendemos na elaboração da proposta}

\textbf{1. Formular a lacuna com precisão mudou o argumento.} Uma hipótese inicial dizia que o PMBOK 8 tratava privacidade apenas de forma implícita. A leitura do guia mostrou menções explícitas a proteção de dados, GDPR e privacidade, inclusive no contexto de inteligência artificial. A formulação defensável passou a ser mais específica: entre os sete domínios e 40 processos oficiais, não há um domínio e um conjunto de processos dedicados a gerir os efeitos do tratamento de dados pessoais sobre indivíduos. Essa revisão evitou construir a extensão sobre uma ausência absoluta que o próprio texto não sustenta.

\textbf{2. Partes interessadas e titulares se sobrepõem, mas não são equivalentes.} Os processos de Partes Interessadas ajudam a identificar pessoas e a organizar comunicação e engajamento. Uma pessoa pode, porém, ter seus dados tratados sem participar do projeto ou conhecer sua existência. O domínio proposto precisa dar visibilidade a esses titulares e às operações sobre os dados, preservando as interfaces com os processos existentes.

\textbf{3. Risco ao projeto não esgota risco à pessoa.} A estrutura do domínio Riscos foi útil para análise, respostas e monitoramento. O NIST Privacy Framework ajudou a distinguir danos possíveis do próprio processamento de dados de incidentes de cibersegurança. Por isso, a proposta associa cada avaliação a pessoas afetadas, finalidade e decisão, sem criar um sistema paralelo que ignore o registro geral de riscos.

\textbf{4. Propor processos exige mais do que nomes e áreas de foco.} O PMBOK 8 descreve cada processo por entradas, ferramentas e técnicas e saídas. Nossa primeira versão trazia entradas e saídas, mas não as técnicas. Acrescentar ferramentas e técnicas com seu fundamento (LGPD, NIST, literatura) tornou os processos verificáveis e mostrou onde a proposta ainda é frágil.

\textbf{5. Um domínio novo só convence se mostrar o que muda nos domínios existentes.} Inicialmente descrevemos apenas interfaces. Ao percorrer os treze processos de Partes Interessadas e Riscos, um a um, ficou claro que a contribuição está nas inclusões e alterações mínimas em cada um (entradas, registros e critérios), mais do que nos oito processos novos isoladamente.

\textbf{6. A literatura sobre gêmeos digitais mudou o foco.} Os estudos examinados tratam de requisitos legais, arquiteturas e controles técnicos. A distinção entre modelo, sombra e gêmeo digital pela direção do fluxo de dados mostrou que o retorno automático ao mundo físico é o ponto em que a privacidade mais exige decisão de gestão. Isso levou a incluir a direção e a frequência dos fluxos no inventário inicial. A observação vale para os estudos que examinamos, não para toda a literatura.

\textbf{7. O escopo exigiu uma decisão editorial explícita.} Desenvolvemos um único domínio, \textbf{Privacidade \ensuremath{\rightarrow} Quem}, e mostramos Sustentabilidade, Inovação e Análise de Negócios somente no mapa de integração. Essa escolha mantém a contribuição revisável e deve permanecer visível na revisão conjunta.

\textbf{8. Um cenário hipotético tornou a proposta testável.} Um gêmeo digital de reciclagem por cooperativas, com mobilização de moradores, triagem e destinação final, permitiu perguntar onde se identificam titulares, como se limita compartilhamento e quem responde pelo destino dos registros. O cenário é uma ilustração dos autores; não foi usado como prova de conformidade de uma cooperativa real.

\subsection*{O que revisaríamos numa próxima versão}

Testaríamos os oito processos com profissionais de projetos, privacidade e cooperativas para eliminar duplicações e calibrar o esforço necessário nos três perfis de contexto. Conferiríamos se cada processo tem entrada disponível, técnica aplicável, decisão atribuída, saída útil e interface clara com Partes Interessadas e Riscos. Testaríamos a hipótese de que maior frequência e resolução dos dados aumentam a possibilidade de identificação em projetos de gêmeo digital com fluxo bidirecional.

\textbf{Limite deste trabalho:} as matrizes e processos são proposta acadêmica documental. Não constituem extensão oficial do PMI, auditoria jurídica ou validação empírica.

\end{document}
